\subsection{Error geometry behind the main analysis}
\label{app:geometry}
Two models with similar individual accuracy can give very different ensemble
gains, depending on where their errors occur and how those errors align.
Classical ensemble analysis makes this dependence explicit
\citep{breiman1996,krogh1994}. For a fixed trained library and population of
inputs, let $C=\E[G(x,y)]$. Because the mixture error is a weighted sum of
individual normalized error fields, its expected squared relative error is
\begin{equation}
 R(w)=\E\left[\frac{\norm{\widehat y_w(x)-y}_Q^2}{\max(\norm y_Q,10^{-8})^2}\right]
 =w^\top Cw.
\label{eq:risk}
\end{equation}
Here $\widehat y_w$ is the combined prediction, $y$ the target, $w$ the model
weights, and $\norm{\cdot}_Q$ the field norm. Expectation $\E$ averages over
input and target pairs, $G(x,y)$ contains their model error products, and $C$
is its mean. The identity computes population squared error from these products.
The diagonal measures individual error and the cross terms measure alignment.
We retain uncentered errors so that shared bias remains part of the risk.
Keeping only a positive diagonal yields the inverse error mixture as the simplex
minimizer before clipping. This is an approximation used to choose weights,
and does not require the actual model errors to be independent.

To separate individual accuracy from complementarity, consider one case and write $r_m=\norm{e_m}_Q$.
If the errors all pointed in the same direction, the squared error of their
weighted sum would be $(\sum_m w_m r_m)^2$.
The difference from the actual mixture is
\begin{equation}
 \left(\sum_m w_m r_m\right)^2-w^\top Gw
 =2\sum_{m<n}w_mw_n\bigl(r_mr_n-\langle e_m,e_n\rangle_Q\bigr)\geq0.
\label{eq:direction}
\end{equation}
Here $r_m$ is the size of error field $e_m$, $G$ is the error product matrix for
the same case, and $m<n$ counts each model pair once. The nonnegative difference
measures how much error alignment matters when the individual error sizes are fixed.
Models can benefit from making errors in different locations or from partially
canceling one another. Negative pairwise correlation is not required.
The inverse error mixture can exploit the same complementarity without fitting cross terms.

\paragraph{The fitting loss can change the preferred weights.}
The effect of the fitting criterion is established in engineering surrogate
ensembles \citep{acar2009}.
The fitted mixture minimizes empirical squared relative error, an estimate of
$R(w)$ in Equation~\eqref{eq:risk}. The reported mean relative $L^2$ error
instead estimates
\begin{equation}
 L(w)=\E\!\left[\sqrt{w^\top G(x,y)w}\right].
\label{eq:reportedrisk}
\end{equation}
Here $L(w)$ is the expected relative $L^2$ error of weights $w$. Taking the
square root for each field before averaging matches the reported metric,
whereas $R(w)$ averages its square.
Jensen's inequality gives $L(w)\leq\sqrt{R(w)}$ for a given $w$ but does not
transfer an excess risk guarantee relative to the minimizer of $L$.
Consider two equally probable scalar cases with normalized errors
$e_A=(0,2)$ and $e_B=(1.1,1.1)$ across cases. Giving A weight $t\in[0,1]$
and B weight $1-t$ yields
\begin{align}
 L(t)&=1.1-0.1t,\\
 R(t)&=1.21-0.22t+1.01t^2.
\end{align}
Here $t$ is model A's weight, and $L(t)$ and $R(t)$ evaluate the unsquared and
squared losses of the same two model mixture. The formulas expose their different optima.
Thus $L$ is minimized at $t=1$, with value one, while $R$ is minimized at
$t=11/101$, where $L=110/101$. Optimizing squared error is approximately
$8.9\%$ worse under the reporting metric in this example. The example also
applies to constant fields and persists with infinitely many fitting cases.

Direct fitting of the reported metric is a convex alternative:
\begin{equation}
 \min_{w\in\Delta_M}\frac{1}{K}\sum_{i=1}^{K}\sqrt{w^\top G_iw}
 =\min_{w\in\Delta_M}\frac{1}{K}\sum_{i=1}^{K}\norm{A_iw}_2,
 \qquad A_i^\top A_i=G_i.
\label{eq:metricfit}
\end{equation}
Here $A_i$ is a matrix factor of the error matrix $G_i$ for fitting example $i$, $\norm{\cdot}_2$
is the Euclidean norm, and $\Delta_M$ enforces nonnegative unit sum weights.
The objective fits the mean relative error over $K$ whole fields directly.
Each summand is a norm of a linear map. Introducing an upper bound $t_i$ for
each field's error, with constraints $\norm{A_iw}_2\leq t_i$, gives a second order cone problem.
This is a standard loss choice, not a new ensemble construction.

\paragraph{What the estimation bound establishes.}
Appendix~\ref{app:proof} shows that entrywise matrix error at most $\epsilon$
and empirical optimization error at most $\delta$ imply excess squared risk
at most $2\epsilon+\delta$ over the best fixed convex mixture for squared
relative error. The deterministic implication requires accurate error
estimation, not independent cases. Independence is needed for the stated
probabilistic sufficient condition, conditional on the frozen library.
That condition concerns whole fields, not individual pixels, and its
assumptions are unverified for this collection. Its numerical bound is weaker
than the trivial bounded risk guarantee at all three reported budgets.
It therefore establishes neither that five fields suffice nor near optimality
under Equation~\eqref{eq:reportedrisk}.

\paragraph{Testing the loss choice while keeping predictions fixed.}
To isolate the fitting objective, we evaluate Equation~\eqref{eq:metricfit} on
all 51 historical partitions at $K=5$, keeping the original nine models,
fitting rows, and evaluation rows. We first reproduce the original selected,
inverse, and fitted error vectors from the stored per case Gram matrices and
verify their input identities. We then fit and save all new weights before
scoring them. The specification was saved before these additional errors were
computed, but the predictions had previously been inspected. This is therefore
a post hoc diagnostic, and the new fit is reported separately from the three main rules.

The epigraph problem uses CVXPY 1.9.2 and Clarabel 0.11.1. All 51 retained solves returned
optimal status with absolute and relative gap tolerances $10^{-10}$ and
feasibility tolerance $10^{-9}$. Scaling by the mean diagonal on the fitting set improves
conditioning without changing the minimizer. Eigenvalues below zero are
clipped only within $10^{-12}$ of spectral scale and every adjustment is
recorded. The most negative relative eigenvalue was approximately
$-1.04\times10^{-18}$. More substantial indefiniteness raises an error.
The largest simplex feasibility residual before nonnegative clipping and
renormalization was $4.15\times10^{-11}$. We also verified the counterexample
above numerically and rejected a deliberately indefinite input matrix.

\begin{table}[t]
\centering\small
\caption{\textbf{Fitting the reported metric directly.} Ratios compare with the
selected model on the original 17 PDE dataset diagnostic subset at five fitting examples.
The fitted mixture uses squared relative error. A win or loss requires
more than a one percent difference after averaging the three partitions for
each dataset. Remaining datasets are within one percent. These are post hoc
results on a fixed library, without independent training replicates.}
\label{tab:losscontrol}
\begin{tabular}{lrrrr}
\toprule
Method & Class ratio & Dataset ratio & Wins & Losses \\
\midrule
Inverse error mixture & 0.882 & 0.924 & 11 & 4 \\
Fitted mixture & 0.876 & 0.918 & 12 & 5 \\
Relative $L^2$ fit & 0.859 & 0.898 & 12 & 5 \\
\bottomrule
\end{tabular}

\end{table}

Direct relative $L^2$ fitting reduces class balanced error by $2.0\%$ and equal
dataset error by $2.1\%$ relative to the original squared error fit
(Table~\ref{tab:losscontrol}). Nine datasets improve by more than one percent,
none worsens by that margin, and eight remain within one percent. Against the
selected model, 12 datasets improve and five worsen. The loss mismatch therefore
has a measurable consequence in these stored predictions. The result motivates
an independent evaluation of this objective choice, without establishing a new
weighting principle or eliminating the limitations of the historical comparison.
